\section{小波与多尺度分解}
\label{sec:08-Numerical-Analysis/08-Fourier-and-Spectral-Methods/06}

一段 \(10\) 秒的心电记录，采样率 \(500\) Hz。上面骑着两样脏东西：呼吸带来的基线漂移，周期在秒的量级；肌电干扰，弥散在几十到上百赫兹。要交付的是干净波形，而且 QRS 那个宽度只有几毫秒的尖峰，位置和幅度都不能动。

先试低通滤波。截止频率往下压，肌电噪声确实小了，QRS 也跟着钝了——尖峰之所以尖，靠的正是几十到上百赫兹那一大片成分，它和噪声挤在同一个频段里，用一把只认频率的刀切不开。再试上一节的\hyperref[sec:08-Numerical-Analysis/08-Fourier-and-Spectral-Methods/05]{``时频局部化''}：给窗口定宽度时，要看清几毫秒的尖峰得用窄窗，可窄窗根本装不下一个周期以秒计的基线漂移。一个信号里同时住着毫秒级和秒级的结构，单一窗宽照顾不了两头。

出路上一节末尾已经点到：让窗宽随频率走。

\subsection{把缩放做成基的组织方式}

放弃``一个窗口配所有频率''的安排，改用一族由同一个母小波（mother wavelet）\(\psi\) 经缩放与平移生成的函数：

\[
\psi_{a,b}(t)=\frac1{\sqrt{\abs a}}\,\psi\!\left(\frac{t-b}{a}\right),\qquad a\neq0,\quad b\in\R .
\]

\(b\) 管平移，也就是何时；\(a\) 管缩放，也就是多宽；因子 \(1/\sqrt{\abs a}\) 让所有小波的 \(L^2\) 范数相同，不同尺度的系数才可比。连续小波变换是往这族函数上投影：

\[
W_f(a,b)=\ip{f}{\psi_{a,b}}=\frac1{\sqrt{\abs a}}\int_{-\infty}^{\infty}f(t)\,\overline{\psi\!\left(\frac{t-b}{a}\right)}\,dt .
\]

\(\abs a\) 小，小波被压窄，对快速变化敏感；\(\abs a\) 大，小波被拉宽，对缓慢轮廓敏感。母小波得是真正的``小''波：振荡、迅速衰减，并且零均值 \(\int\psi(t)\,dt=0\)——否则它对信号的常数部分也有响应，就谈不上只抓变化了。

关键在于窗宽和频率是绑在一起变的。尺度 \(a\) 处的小波主要响应 \(1/a\) 附近的频率，同时它在时间上的支撑也正比于 \(a\)。于是频率越高，时间窗自动越窄；频率越低，时间窗自动越宽。上一节要靠人手挑的那个参数，在这里成了基的组织方式本身。

\begin{center}
\begin{tikzpicture}[scale=1.0,font=\small,>=stealth]
  \draw[->,black!55] (-0.1,0) -- (4.7,0) node[right,font=\footnotesize] {\(t\)};
  \draw[->,black!55] (0,-0.1) -- (0,4.5) node[above,font=\footnotesize] {\(\omega\)};
  \foreach \i in {0,1}{
    \draw[black!55,fill=black!6] (\i*2.0,0) rectangle (\i*2.0+2.0,0.5);
  }
  \foreach \i in {0,1}{
    \draw[orange!70!black,fill=orange!8] (\i*2.0,0.5) rectangle (\i*2.0+2.0,1.0);
  }
  \foreach \i in {0,1,2,3}{
    \draw[orange!70!black,fill=orange!8] (\i*1.0,1.0) rectangle (\i*1.0+1.0,2.0);
  }
  \foreach \i in {0,1,2,3,4,5,6,7}{
    \draw[orange!70!black,fill=orange!8] (\i*0.5,2.0) rectangle (\i*0.5+0.5,4.0);
  }
  \node[right,font=\footnotesize,black!70] at (4.15,0.25) {近似};
  \node[right,font=\footnotesize,black!70] at (4.15,0.75) {粗尺度细节};
  \node[right,font=\footnotesize,black!70] at (4.15,3.0) {细尺度细节};
  \draw[<->,black!70] (0.05,4.28) -- (0.45,4.28);
  \node[above,font=\footnotesize,black!70] at (0.25,4.3) {窄};
  \draw[<->,black!70] (2.05,4.28) -- (3.95,4.28);
  \node[above,font=\footnotesize,black!70] at (3.0,4.3) {宽};
  \node[font=\footnotesize,black!70] at (2.0,-0.55) {每块面积仍然相同，长宽比随频率翻转};
\end{tikzpicture}

\emph{回到上一节那张图：那里的格子形状全同，横竖比例一选定就管住整个平面。这里高频那几层的格子窄而高，低频那几层宽而矮，面积却和上一节一样压不下去——变的是分配方式，不是总量。}
\end{center}

图里每块砖的面积都没变，不确定性原理的下界一分没让。变的只是同一份预算在不同频段上的花法：高频处买时间精度，低频处买频率精度。心电那个例子于是有了着落——尖峰落在最上面几层的少数几块砖里，基线漂移落在最下面一层，两者在坐标上就分开了。

还要提醒一句：小波占的是一段频带而不是单一频率，``尺度''只是频率的粗糙代理。两者不能当作同一个量互相换算。

\subsection{Haar 小波是这套想法最短的完整例子}

在 \([0,1)\) 上，令 \(V_j\) 为在 \(2^j\) 个等长二分子区间 \([k2^{-j},(k+1)2^{-j})\) 上分别取常数的平方可积函数全体，即``分辨率 \(j\) 下的分段常数近似''。格子每加密一倍，能表示的函数就更多，于是 \(V_0\subset V_1\subset V_2\subset\cdots\)。\(V_j\) 的天然正交基由尺度函数（scaling function）\(\phi=\ind_{[0,1)}\) 缩放平移给出：\(\phi_{j,k}(t)=2^{j/2}\phi(2^jt-k)\)，其中 \(\ind\) 是指示函数，因子 \(2^{j/2}\) 保持单位范数。

要问的是：从 \(V_j\) 升到 \(V_{j+1}\)，新增的自由度是什么？一个区间上的常数拆成左右两半上的两个常数，多出来的信息恰好是左半减右半这一个数。照着这句话定义 Haar 母小波

\[
\psi(t)=
\begin{cases}
1, & 0\le t<1/2,\\
-1, & 1/2\le t<1,
\end{cases}
\]

并记 \(\psi_{j,k}(t)=2^{j/2}\psi(2^jt-k)\)。每个 \(\psi_{j,k}\) 与所有 \(\phi_{j,\ell}\) 正交，彼此也正交，于是

\[
V_{j+1}=V_j\oplus W_j,\qquad W_j=\operatorname{span}\set{\psi_{j,k}}_{k=0}^{2^j-1}.
\]

\(W_j\) 叫细节空间，它就是相邻两级分辨率之差。近似归 \(V_j\)，差出来的细节归 \(W_j\)，两边正交，没有重复计数。

\subsection{手算一层，看细节系数在哪里鼓起来}

取 8 个样本

\[
x=(8,4,6,2,7,1,3,5),
\]

看作 \(V_3\) 中某个函数在 8 个区间上的取值。相邻两两配对，每对算均值与半差：

\[
a_m=\frac{x_{2m}+x_{2m+1}}2,\qquad d_m=\frac{x_{2m}-x_{2m+1}}2 .
\]

逐对代进去：\((8,4)\mapsto(6,2)\)，\((6,2)\mapsto(4,2)\)，\((7,1)\mapsto(4,3)\)，\((3,5)\mapsto(4,-1)\)。得近似系数 \(a=(6,4,4,4)\) 与细节系数 \(d=(2,2,3,-1)\)。8 个数变成 \(4+4\) 个数，一半记概貌，一半记变化。重构是平凡的：\(x_{2m}=a_m+d_m\)，\(x_{2m+1}=a_m-d_m\)，代回原样取回 \((8,4,6,2,7,1,3,5)\)，一位信息都没丢。

这就是 \(V_3=V_2\oplus W_2\) 的坐标版本。细节系数值得盯一眼：第三个是 \(3\)，明显大过其余——原序列里 \((7,1)\) 那一对落差最大，局部变化在细节系数上直接鼓起一个包。心电里的 QRS 尖峰在细尺度细节上做的就是这件事。

归一化要交代一句。上面的均值与半差约定让重构最简，但每层的能量关系是 \(\norm{x}^2=2(\norm{a}^2+\norm{d}^2)\)。把两个 \(1/2\) 都换成 \(1/\sqrt2\)，每层就成了严格正交的坐标变换，能量精确守恒，对应前面那组单位范数的基。这与 DFT 里 \(1/N\) 因子放在哪一侧是同类问题：比较任何数值之前先对约定。

\subsection{把分解叠成塔}

Haar 例子里起作用的是一套可以抽出来的结构，具体公式只是它的一个实现。多分辨率分析（multiresolution analysis, MRA）要求一串闭子空间 \(V_j\) 满足四条：嵌套，\(V_j\subset V_{j+1}\)，分辨率只增不减；完备，所有 \(V_j\) 的并在 \(L^2\) 中稠密，交则退化（在 \(\R\) 上只剩零函数，在 \([0,1)\) 上只剩常数，上例最后剩下的那个总均值正是它）；缩放律，\(f(t)\in V_j\iff f(2t)\in V_{j+1}\)，相邻两级只差二倍缩放；以及存在尺度函数 \(\phi\)，其整数平移构成 \(V_0\) 的正交基。

有了这四条，细节空间 \(W_j\) 就定义成 \(V_j\) 在 \(V_{j+1}\) 里的正交补，而且分解可以对近似部分反复迭代：

\[
V_J=V_0\oplus W_0\oplus W_1\oplus\cdots\oplus W_{J-1}.
\]

前面那 8 个数继续往下走就是一座塔。近似 \((6,4,4,4)\) 再配对，得近似 \((5,4)\) 与细节 \((1,0)\)；\((5,4)\) 再来一次，得总均值 \(4.5\) 与细节 \(0.5\)。原信号最终写成总均值 \(4.5\)，加上 \(W_0\) 细节 \((0.5)\)、\(W_1\) 细节 \((1,0)\)、\(W_2\) 细节 \((2,2,3,-1)\)，共 \(1+1+2+4=8\) 个系数，依旧无损。粗尺度系数少而管整体，细尺度系数多而管局部，同一个信号被远观和近看同时记录了下来。这也正是图像金字塔和各类多尺度特征在做的事：一层粗糙轮廓加若干层逐渐变细的差异。

Haar 不是唯一选择。Daubechies 小波是它的光滑推广，用更长的滤波器换来更光滑的基函数和更多消失矩（vanishing moments）。\(p\) 阶消失矩意味着小波对次数低于 \(p\) 的多项式成分响应为零，因此光滑段上的细节系数自动接近零，稀疏性更好；Haar 是其中最简的一员，只有 \(1\) 阶。代价是支撑变长、边界处理更麻烦，而且可以证明，紧支撑的正交实小波里除 Haar 外没有对称的。

计算上，配对求均值与半差就是两个滤波器：低通 \(h\) 负责平均，高通 \(g\) 负责差分，各接一次二倍下采样，再对低通输出递归。这是 Mallat 的快速小波变换，每层线性工作量，整座塔加起来 \(O(N)\)。\hyperref[sec:08-Numerical-Analysis/08-Fourier-and-Spectral-Methods/03]{``DFT 是变换，FFT 是算法''}那条分界在这里原样重现：小波基是数学对象，滤波器组塔是算它的算法，而这个算法比 FFT 的 \(O(N\log N)\) 还便宜一档。

\subsection{稀疏在哪里，小波就在哪里有用}

分段光滑的信号在小波坐标下稀疏：能量集中在少数几个大系数里，其余系数贴着零。白噪声没有这个性质，它在正交变换下仍是白噪声，均匀摊到所有系数上。两件事一叠，去噪的做法就自明了——把小系数置零（软阈值或硬阈值），再重构。噪声被大面积削掉，而承载跳变边沿的那几个大系数原封不动。开头那段心电用这个办法处理，QRS 不会被削平，因为它的能量不是摊在整个高频段里，而是攒在少数几块砖上。低通滤波做不到这一点，它只认频率不认位置。

压缩是同一件事换个目的：光滑区域的细节系数本就近零，只需存下少数大系数。JPEG2000 用离散小波变换替下 JPEG 的分块离散余弦变换，看中的正是这份稀疏性，顺带避开了分块带来的块效应。突变检测则利用尺度间的一致性：真正的局部奇异（尖峰、断点、边沿）会在连续好几个尺度的细节系数上都留下大值，而噪声制造的极值随尺度变粗迅速消失，两者可以分开。

边界也清楚。若信号是平稳谐波，整段由固定频率的正弦叠成，Fourier 基就是它的本征坐标，几个系数解决问题；换成小波反而把能量摊到许多尺度上，既费系数又难读。小波的地盘是分段光滑、含局部奇异、统计性质随时间变化的那一类信号。即便在这一类里面，选 Haar 还是 Daubechies、选几阶消失矩，也都是建模决定而非技术细节：消失矩阶数决定了什么样的``光滑''会被判为无信息，对形态不同的信号，判断并不一致。没有哪一族基对所有信号都稀疏，这是没有免费午餐在信号表示里的形态。

回头看整个单元，六节做的是同一件事的六个层次：把函数换一组坐标，让原本纠缠的性质变成坐标上的读数。Fourier 坐标让微分成为乘法、卷积成为逐点相乘、光滑性成为系数衰减；采样与 DFT 交代了这组坐标离散化之后会漏掉什么；谱方法把系数衰减兑换成计算精度；加窗与小波则在这组坐标丢掉位置信息之后，一步步把位置找回来。

留了一笔账没结。谱方法那节说，解析函数的系数按 \(e^{-a\abs{k}}\) 衰减，精度全看这个 \(a\)；可 \(a\) 到底是多少，实轴上无从计算——两个一样光滑、一样有界的函数，衰减率可以差出好几倍。下一章\hyperref[ch:069]{``复分析与留数：奇点决定函数的深层结构''}把坐标继续往外挪一步，挪进复平面。那里 \(a\) 有一个干脆的答案：它是函数最近的那个奇点到实轴的距离。
